\subsection{代数的深度}

引理 2.--- 设 \(\rho: A \to B\) 为 Noether 局部环之间的局部同态，N 为有限生成 B-模，y 为 \(\mathfrak{m}_{\mathrm{B}}\) 的元素。下列两个条件等价：

\begin{enumerate}
\def\labelenumi{(\roman{enumi})}
\item
  A-模 N/yN 平坦且乘以 y 的映射 \(y_{N}\) 单射；
\item
  A-模 N 平坦且乘以 y 的映射 \(y_{\kappa_{A}\otimes N}\) 单射。
\end{enumerate}

当这些条件满足时，对每个 A-模 M，乘以 y 的映射 \(y_{\mathrm{M} \otimes_{\Lambda}\mathrm{N}}\) 都单射。

设 (i) 的假设成立，并同时证明 (ii) 及最后一个断言。设 M 为 A-模。由于 A-模 N/yN 平坦，由正合序列 \(0 \rightarrow N \xrightarrow{y_{N}} N \longrightarrow N/yN \rightarrow 0\) 推出正合序列

\[
0 \rightarrow \mathrm{M} \otimes_ {\mathrm{A}} \mathrm{N} \xrightarrow {u} \mathrm{M} \otimes_ {\mathrm{A}} \mathrm{N} \longrightarrow \mathrm{M} \otimes_ {\mathrm{A}} (\mathrm{N} / y \mathrm{N}) \rightarrow 0
\]

\[
0 \to \operatorname{Tor} _ {1} ^ {\mathrm{A}} (\mathrm{M}, \mathrm{N}) \stackrel {v} {\longrightarrow} \operatorname{Tor} _ {1} ^ {\mathrm{A}} (\mathrm{M}, \mathrm{N}) \to 0
\]

其中 \(u = 1_{\mathrm{M}} \otimes y_{\mathrm{N}}\) ， \(v = \operatorname{Tor}_1^{\mathrm{A}}(1_{\mathrm{M}}, y_{\mathrm{N}})\) ；由此得出，乘以 \(y\) 的映射在 \(\mathrm{M} \otimes_{\mathrm{A}} \mathrm{N}\) 中是单射的，且在 \(\operatorname{Tor}_1^{\mathrm{A}}(\mathrm{M}, \mathrm{N})\) 中是双射的。再设 A-模 \(\mathrm{M}\) 是有限型的；则 B-模 \(\operatorname{Tor}_1^{\mathrm{A}}(\mathrm{M}, \mathrm{N})\) 是有限型的 (A, X, 第 107 页, 命题 6)，由于 \(y\) 属于 \(\mathfrak{m}_{\mathrm{B}}\) ，它由 Nakayama 引理为零，这蕴含 A-模 \(\mathrm{N}\) 平坦 (A, X, 第 74 页, 定理 2)。

(ii)⇒(i)：设 的假设成立。考察 B-模的下列两个正合序列

\begin{enumerate}
\def\labelenumi{(\arabic{enumi})}
\tightlist
\item
\end{enumerate}

\[
0 \to \mathrm{Ker} (y _ {\mathrm{N}}) \longrightarrow \mathrm{N} \stackrel {p} {\longrightarrow} \mathrm{Im} (y _ {\mathrm{N}}) \to 0\tag{1}
\]

\[
0 \to \mathrm{Im} (y _ {\mathrm{N}}) \stackrel {i} {\longrightarrow} \mathrm{N} \longrightarrow \mathrm{N} / y \mathrm{N} \to 0,\tag{2}
\]

其中 \(p\) 与 \(i\) 是典范同态。由此得出，同态 \(1 \otimes p : \kappa_{\mathrm{A}} \otimes_{\Lambda} \mathrm{N} \longrightarrow \kappa_{\mathrm{A}} \otimes_{\mathrm{A}} \mathrm{Im}(y_{\mathrm{N}})\) 是满的，且（由于 N 平坦）同态 \(1 \otimes i : \kappa_{\mathrm{A}} \otimes_{\mathrm{A}} \mathrm{Im}(y_{\mathrm{N}}) \longrightarrow \kappa_{\mathrm{A}} \otimes_{\Lambda} \mathrm{N}\) 的核同构于 \(\operatorname{Tor}_1^{\mathrm{A}}(\kappa_{\mathrm{A}}, \mathrm{N}/y\mathrm{N})\) 。而映射 \((1 \otimes i) \circ (1 \otimes p)\) ——它等于 \(y_{\kappa_{\mathrm{A}} \otimes_{\mathrm{A}} \mathrm{N}}\) ——按假设是单射；由此得出 \(1 \otimes p\) 双射且 \(1 \otimes i\) 单射，从而有 \(\operatorname{Tor}_1^{\mathrm{A}}(\kappa_{\mathrm{A}}, \mathrm{N}/y\mathrm{N}) = 0\) 。于是得出 A-模 \(\mathrm{N}/y\mathrm{N}\) 平坦 (III, § 5, 第 2 节, 定理 1 及第 4 节, 命题 2)。

由于 N 与 N/yN 在 A 上平坦，\(\mathrm{Im}(y_{\mathrm{N}})\) 亦然（正合序列 (2)）。于是由正合序列 (1) 推出，\(\kappa_{A} \otimes_{A} \mathrm{Ker}(y_{\mathrm{N}})\) 同构于 \(1 \otimes p\) 的核；故它为零；从而乘以 y 的映射 \(y_{N}\) 由 Nakayama 引理是单射的。

命题 10. 设 \(\rho: \Lambda \to B\) 为 Noether 局部环之间的局部同态，N 为有限生成 B-模，\(y = (y_1, \ldots, y_s)\) 为 \(\mathfrak{m}_{\mathrm{B}}\) 中元素组成的序列。设 \(\mathfrak{y}\) 为由该序列生成的 B 的理想。下列条件等价：

\begin{enumerate}
\def\labelenumi{(\roman{enumi})}
\item
  A-模 N/\(\mathfrak{y}\)N 平坦且序列 \(y\) 是 N-正则的；
\item
  A-模 N 平坦且序列 \(y\) 是 \((\kappa_{A} \otimes_{A} N)\) -正则的。
\end{enumerate}

当这些条件满足时，对每个 A-模 M，序列 \(y\) 是 \(M\otimes_{\mathrm{A}}N\) -正则的。

我们对 s 作归纳来证明 与 的等价性。s = 0 的情形是显然的，设 \(s \geqslant 1\) 。用 \(y'\) 表示序列 \((y_{1}, \ldots, y_{s-1})\) ，用 \(\mathfrak{y}'\) 表示由它生成的 B 的理想。把引理 2 应用于 B-模 N/ \(y'\) N 及 B 的元素 \(y_{s}\) ，条件 等价于

(i') A-模 N/ \(\mathfrak{y}'\) N 平坦，序列 \(y'\) 是 N-正则的，且乘以 \(y_s\) 的映射在 \(\kappa_{\mathrm{A}} \otimes_{\mathrm{A}} (\mathrm{N}/\mathfrak{y}'\mathrm{N}) = (\kappa_{\mathrm{A}} \otimes_{\mathrm{A}} \mathrm{N})/\mathfrak{y}'(\kappa_{\mathrm{A}} \otimes_{\mathrm{A}} \mathrm{N})\) 中是单射的。

由归纳假设，该条件等价于。

最后一个断言类似地由引理 2 的最后一个断言对 \(s\) 作归纳得出。

命题 11. 设 \(\rho: A \to B\) 为 Noether 局部环之间的局部同态，M 为有限生成 A-模，N 为有限生成 B-模；我们假设 A-模 N 平坦。

\begin{enumerate}
\def\labelenumi{\alph{enumi})}
\item
  设 \((x_{1},\ldots,x_{r})\) 为 \(m_{A}\) 中元素组成的、对 A-模 M 正则的序列，\((y_{1},\ldots,y_{s})\) 为 \(m_{B}\) 中元素组成的、对 B-模 \(\kappa_{A}\otimes_{A}N\) 正则的序列；则序列 \((y_{1},\ldots,y_{s},\rho(x_{1}),\ldots,\rho(x_{r}))\) ——由 \(m_{B}\) 中元素组成——对 B-模 \(M\otimes_{A}N\) 正则。
\item
  我们有等式
\end{enumerate}

\[
\mathrm{prof} _ {\mathrm{B}} (\mathrm{M} \otimes_ {\mathrm{A}} \mathrm{N}) = \mathrm{prof} _ {\mathrm{A}} (\mathrm{M}) + \mathrm{prof} _ {\mathrm{B}} (\kappa_ {\mathrm{A}} \otimes_ {\mathrm{A}} \mathrm{N}).
\]

设 \(\mathfrak{x}\) 为由序列 \(\mathbf{x}\) 生成的 A 的理想，\(\mathfrak{y}\) 为由 \(y\) 生成的 B 的理想。由命题 10，序列 \(y\) 是 \(M \otimes_A N\) -正则的，且 \(N/\mathfrak{y}N\) 在 \(A\) 上平坦，于是序列 \(\boldsymbol{\rho}(\mathbf{x}) = (\boldsymbol{\rho}(x_1), \ldots, \boldsymbol{\rho}(x_r))\) 对 \(M \otimes_A (N/\mathfrak{y}N) = (M \otimes_A N)/\mathfrak{y}(M \otimes_A N)\) 正则。这就证明了 a)。

为证明 b)，我们可设 M 与 N 非零。由 Nakayama 引理，\(\kappa_{\mathrm{A}} \otimes_{\mathrm{A}} \mathrm{N}\) 也非零，故 \(\operatorname{prof}_{\mathrm{A}}(\mathrm{M})\) 与 \(\operatorname{prof}_{\mathrm{B}}(\kappa_{\mathrm{A}} \otimes_{\mathrm{A}} \mathrm{N})\) 有限（第 4 节, 定理 2 的推论 2）。取极大的正则序列 x 与 y；于是有 \(r = \operatorname{prof}_{\mathrm{A}}(\mathrm{M})\) ， \(s = \operatorname{prof}_{\mathrm{B}}(\kappa_{\mathrm{A}} \otimes_{\mathrm{A}} \mathrm{N})\) ，且存在单 A-线性映射 \(u: \kappa_{\mathrm{A}} \to \mathrm{M} / x \mathrm{M}\) 与单 B-线性映射 \(v: \kappa_{\mathrm{B}} \to \kappa_{\mathrm{A}} \otimes_{\mathrm{A}} (\mathrm{N} / \mathfrak{y} \mathrm{N})\) 。由于 N/ \(\mathfrak{y}\) N 在 A 上平坦，B-线性映射 \((u \otimes 1_{\mathrm{N} / \mathfrak{y} \mathrm{N}}) \circ v\) ——它把 \(\kappa_{\mathrm{B}}\) 映入 (M/xM) \(\otimes_{\mathrm{A}} (\mathrm{N} / \mathfrak{y} \mathrm{N}) = (\mathrm{M} \otimes_{\mathrm{A}} \mathrm{N}) / (\rho(x) + \mathfrak{y})(\mathrm{M} \otimes_{\mathrm{A}} \mathrm{N})\) ——是单射的。这蕴含等式 \(\operatorname{prof}_{\mathrm{B}}(\mathrm{M} \otimes_{\mathrm{A}} \mathrm{N}) = r + s\) （同上引文）。

注. -- 回忆在此前诸假设下，我们有

\[
\dim_ {B} (M \otimes_ {A} N) = \dim_ {A} (M) + \dim_ {B} (\kappa_ {A} \otimes_ {A} N)
\]

(VIII, § 3, 第 4 节, 命题 7)。

推论.--- 设 \(\rho: A \to B\) 为 Noether 局部环之间的使 B 成为平坦 A-模的局部同态。我们有

\[
\begin{array}{l} \operatorname{prof} (\mathrm{B}) = \operatorname{prof} (\mathrm{A}) + \operatorname{prof} \left(\kappa_ {\Lambda} \otimes_ {\Lambda} \mathrm{B}\right), \\ \dim (\mathrm{B}) = \dim (\mathrm{A}) + \dim \left(\kappa_ {\mathrm{A}} \otimes_ {\mathrm{A}} \mathrm{B}\right). \end{array}
\]

事实上，由命题 4 的推论（相应地，由 VIII, § 1, 第 4 节），B-模 \(\kappa_{\mathrm{A}}\otimes_{\mathrm{A}}\mathrm{B}\) 的深度（相应地，维数）等于环 \(\kappa_{\mathrm{A}}\otimes_{\mathrm{A}}\mathrm{B}\) 的深度（相应地，维数）。

